\chapter{测度、可测函数与 Lebesgue 积分}
\label{ch:measure-integration}

\begin{introduction}
  \item 先修：第 2 章的极限、完备性、函数列逐点与一致收敛。
  \item 学习产出：能写出概率模型的可测结构，并判断极限与积分交换缺少哪条假设。
  \item 研究连接：随机变量、期望、损失函数和连续时间模型都建立在同一测度积分语言上。
\end{introduction}

Riemann 积分从区间分割出发；概率与现代统计更需要从“哪些事件可测、事件有多大、函数如何由简单对象逼近”出发。测度论的作用不是把积分写得更抽象，而是让零概率例外、函数极限和期望交换具有可判定前提。

\section{$\sigma$ 代数：先规定哪些集合可被提问}

集合 $\Omega$ 上的集合族 $\mathcal{F}$ 若包含 $\Omega$，对补集封闭，并对可数并封闭，就称为 $\sigma$ 代数：
\[
\Omega\in\mathcal{F},\qquad
A\in\mathcal{F}\Rightarrow A^c\in\mathcal{F},\qquad
A_n\in\mathcal{F}\Rightarrow\bigcup_{n=1}^{\infty}A_n\in\mathcal{F}.
\]
由 De Morgan 律，它也对可数交封闭。$\mathbb{R}$ 上包含所有开集的最小 $\sigma$ 代数称为 Borel $\sigma$ 代数。它已经覆盖通常数据事件，但不等于幂集；Vitali 集说明，若坚持平移不变、可数可加且区间长度正确，就不能让每个子集都成为 Lebesgue 可测集。

\begin{note}
数据表中的“所有能写出的筛选条件”不是自动可测。实际模型应先声明样本空间、事件 $\sigma$ 代数以及数据生成变量的可测性。
\end{note}

\section{测度：从事件到可数可加的大小}

测度 $\mu:\mathcal{F}\to[0,\infty]$ 满足 $\mu(\varnothing)=0$，且对两两不交的 $A_n$ 有
\[
\mu\!\left(\bigcup_{n=1}^{\infty}A_n\right)
=\sum_{n=1}^{\infty}\mu(A_n).
\]
三元组 $(\Omega,\mathcal{F},\mu)$ 是测度空间。Lebesgue 测度给区间赋予长度；若 $\mu(\Omega)=1$，它就是概率测度。可数可加性推出单调性与从下连续性：若 $A_n\uparrow A$，则 $\mu(A_n)\uparrow\mu(A)$。

\section{可测函数：用逆像连接数值与事件}

函数 $f:(\Omega,\mathcal{F})\to(\mathbb{R},\mathcal{B})$ 可测，是指每个 Borel 集 $B$ 的逆像 $f^{-1}(B)$ 都属于 $\mathcal{F}$。等价地，只需检查所有阈值事件 $\{\omega:f(\omega)>t\}$。概率空间上的实值可测函数就是随机变量。

逆像而非像集出现在定义中，因为事件必须留在 $\Omega$ 上。连续函数在 Borel $\sigma$ 代数下可测；可测函数的逐点极限仍可测。这使估计量极限仍是合法随机对象。

\section{简单函数先建立积分账本}

非负简单函数只有有限多个取值，可写成
\[
\phi=\sum_{k=1}^{K}a_k\mathbf{1}_{A_k},
\qquad a_k\ge0,\quad A_k\in\mathcal{F},
\]
其中 $\mathbf{1}_{A}$ 是示性函数。若 $A_k$ 两两不交，则定义
\[
\int\phi\,d\mu=\sum_{k=1}^{K}a_k\mu(A_k).
\]
对任意非负可测函数 $f$，Lebesgue 积分由所有不超过它的简单函数从下逼近：
\[
\int f\,d\mu
=\sup\left\{\int\phi\,d\mu:0\le\phi\le f,\ \phi\text{ 为简单函数}\right\}.
\]

\subsection{固定账本：逼近 $f(x)=x$}

在 $[0,1]$ 上令 $\mu$ 为 Lebesgue 测度，并取
\[
\phi_m(x)=\frac{\lfloor2^m x\rfloor}{2^m}\quad(0\le x<1),
\qquad \phi_m(1)=1.
\]
则 $0\le\phi_m\uparrow f$，且
\[
\int_0^1\phi_m(x)\,dx
=\sum_{k=0}^{2^m-1}\frac{k}{2^m}\frac1{2^m}
=\frac12-\frac1{2^{m+1}}\uparrow\frac12.
\]
因此 $m=2,4,8$ 时的独立结果分别为 $0.375$、$0.46875$ 和 $0.498046875$。

\section{符号函数与可积性}

对实值可测函数，写 $f=f^+-f^-$，其中 $f^+=\max(f,0)$、$f^-=\max(-f,0)$。若两部分积分至少有一个有限，扩展积分有定义；若
\[
\int|f|\,d\mu<\infty,
\]
则称 $f$ 可积，此时 $\int f\,d\mu=\int f^+\,d\mu-\int f^-\,d\mu$ 是有限数。不能把 $\infty-\infty$ 当成零；这也是重尾模型需要先检查绝对可积性的原因。

\section{何时可以交换极限与积分}

三条规则承担不同任务：Fatou 引理给非负函数列的下界；单调收敛定理允许 $0\le f_n\uparrow f$ 时交换；支配收敛定理要求 $f_n\to f$ 几乎处处，且存在可积 $g$ 使 $|f_n|\le g$，此时
\[
\lim_{n\to\infty}\int f_n\,d\mu=\int f\,d\mu.
\]
“几乎处处”表示失败集合的测度为零，不表示有限样本中从未出现。

\subsection{失败见证：质量向零点集中}

在 $[0,1]$ 上令
\[
f_n(x)=n\mathbf{1}_{(0,1/n]}(x).
\]
对每个固定 $x$，都有 $f_n(x)\to0$；但
\[
\int_0^1f_n(x)\,dx=n\cdot\frac1n=1,
\qquad
\int_0^1\lim_{n\to\infty}f_n(x)\,dx=0.
\]
所以极限与积分不能交换。该函数列没有共同的可积支配函数：峰值不断增高并向零点集中。缺失的不是更多网格，而是支配条件。

\section{独立计算与数值边界}

Notebook 用有限求和复现三个简单函数积分，并核对尖峰在 $n=10,100$ 时积分仍为 1、固定点 $x=0.2$ 的函数值已为零：
\begin{lstlisting}[language=bash]
uv run python tools/check_learning_unit.py \
  --manifest curriculum/manifest.json --unit upper.ch03
\end{lstlisting}
离散网格若没有采到 $(0,1/n]$，甚至会把尖峰积分误报为零；解析区间宽度才是本例 oracle。

\section*{章末练习}
\begingroup\small
\subsection*{口述概念}
\begin{enumerate}[nosep,leftmargin=*]
  \item 解释 $\sigma$ 代数、测度、可测函数、简单函数与 Lebesgue 积分的依赖顺序；为什么定义可测函数时使用逆像？
  \item 区分 Borel 可测集、Lebesgue 可测集、零测集与“几乎处处”。为什么 a.e. 结论不等于样本中永不失败？
  \item 比较 Fatou 引理、单调收敛定理和支配收敛定理的假设与结论，并解释共同支配函数为什么必须可积。
\end{enumerate}

\subsection*{笔试推导}
\begin{enumerate}[nosep,leftmargin=*]
  \item 从可数可加性证明测度的单调性、从下连续性；再证明从上连续性并指出 $\mu(A_1)<\infty$ 不能删除。
  \item 推导固定简单函数 $\phi_m$ 的积分公式 $1/2-2^{-(m+1)}$，并用单调收敛定理得到 $\int_0^1x\,dx=1/2$。
  \item 用 Fatou 引理分别作用于 $g+f_n$ 与 $g-f_n$，完成支配收敛定理的积分夹逼证明。
\end{enumerate}

\subsection*{数值编程}
\begin{enumerate}[leftmargin=*]
  \item 计算 $m=2,4,8$ 时逼近 $x$ 的二进台阶积分。
  \item 改变尖峰 $n\mathbf1_{(0,1/n]}$ 的 $n$，比较面积与固定点值。
  \item 比较两种 Pareto 模型的截断期望，并解释能否用 DCT。
\end{enumerate}

\subsection*{研究判断}
\begin{enumerate}[nosep,leftmargin=*]
  \item 研究员看到损失 $L_n(\omega)\to L(\omega)$ a.e.，便宣布期望收敛。列出可以补足结论的两条不同路线及各自证据。
  \item 一组收益数据只有有限个缺失值。能否直接把缺失集合称为零测集并忽略？分别从经验数据与概率模型层评价。
  \item 模拟程序在越来越密的网格上给出稳定积分。设计一份比“再加网格”更强的可积性与极限交换审计。
\end{enumerate}
\endgroup

\subsection*{基础衔接：独立计算}
设 $X$ 的密度为 $2x^{-3}$（$x\geq1$）。求 $\mathbb EX$ 与 $\mathbb EX^2$，判断能否用 $X^2$ 支配函数列并调用 DCT。

\subsection*{分级提示}
\begingroup\small
\begin{enumerate}[nosep,leftmargin=*]
  \item \textbf{口述概念：}先有事件语言，才能赋大小；函数用阈值逆像回到事件；积分先在有限值函数上定义。Borel 与完备化不同，a.e. 是模型中的测度陈述。三条收敛工具分别提供单侧界、单调等式与支配等式。
  \item \textbf{笔试推导：}把递增集合拆成不交增量；从上连续性改看 $A_1\setminus A_n$。简单函数每级高度乘宽度。DCT 对 $g\pm f_n$ 使用 Fatou 后分别控制 $\liminf$ 与 $\limsup$。
  \item \textbf{数值编程：}expected 来自等差和与高度乘宽度；坏输入应从复制的 JSON 进入公开 CLI。固定网格的采样步长若大于峰宽，可能一个正值都看不到。
  \item \textbf{研究判断：}逐点或 a.e. 收敛不控制质量；寻找非负单调性或共同 $L^1$ 支配。有限数据中的“缺失几个”不是自动的概率零测证明；报告测度空间、可测性、尾部界与独立解析 oracle。
\end{enumerate}
\endgroup


\section*{本章小结与 Capstone 连接}

$\sigma$ 代数规定合法事件，测度给事件大小，可测函数把数值问题拉回事件，简单函数构造积分，收敛定理则规定极限交换门禁。上册 Capstone 将据此审计期望、损失与模拟估计中的可测性、可积性和支配条件。
